\documentclass[a4paper,12pt]{scrreprt}
\usepackage{titlesec}
\titleclass{\chapter}{straight}
\titleformat{\chapter}
    {\normalfont\huge\bfseries}{\thechapter\ }{0pt}{\Huge}
\titlespacing*{\chapter}{0pt}{0pt}{20pt}
\usepackage{geometry} %% Used for changing geometry of the page
\usepackage[portuguese]{babel} %% Portuguese language
\usepackage[T1]{fontenc} %% Font setup (pdfLaTeX compatible)
\usepackage[utf8]{inputenc}
\usepackage{helvet}
\renewcommand{\familydefault}{\sfdefault} % Arial-like font
\usepackage{xcolor,colortbl} %% Colors and tables
\usepackage{arydshln}
\usepackage[intoc]{nomencl} %% Acronyms
\makenomenclature
\usepackage{graphicx} %% Images
\graphicspath{{images/}}
\usepackage{hyperref} %% Hyperlinks
\hypersetup{
    colorlinks,
    citecolor=black,
    filecolor=black,
    linkcolor=black,
    urlcolor=black
}
\usepackage{amsmath} %% Math
\usepackage[some]{background} %% Background (for cover)
\usepackage{tikz} %% TikZ graphics
\usetikzlibrary{calc}
\usepackage{listings} %% Code listings
\usepackage{float} %% Better figure placement

%% Configuração do listings para código C
\lstset{
    language=C,
    basicstyle=\ttfamily\footnotesize,
    keywordstyle=\color{blue},
    commentstyle=\color{gray},
    stringstyle=\color{red},
    numbers=left,
    numberstyle=\tiny\color{gray},
    stepnumber=1,
    numbersep=5pt,
    backgroundcolor=\color{white},
    showspaces=false,
    showstringspaces=false,
    frame=single,
    breaklines=true,
    captionpos=b
}

%% ======================= %%
%% DOCUMENT
\begin{document}
%% Blue cover color
\definecolor{titlepagecolor}{RGB}{120,55,59}

%==========================================================================
% COLORED BAR ON THE LEFT SIDE
%==========================================================================

\backgroundsetup{
    scale=1, 
    angle=0, 
    opacity=1,
    contents={
        \begin{tikzpicture}[remember picture,overlay]
            \path [fill=titlepagecolor] 
                (current page.north west) -- ($(current page.north west) + (5,0)$)
                -- ($(current page.south west) + (5,0)$)-- (current page.south west); 
            \node[color=white] at ($(current page.south west) + (3,4)$) {\bfseries {\fontsize{120}{60} \textsf{LI}}};
            \node[color=titlepagecolor] at ($(current page.south west) + (5.8,4)$) {\bfseries {\fontsize{120}{60} \textsf{ 3}}};
        \end{tikzpicture}
    }
}

%==========================================================================
% TITLE PAGE INFO
%==========================================================================

%% Changes values in this field to show information in the cover and back cover about your team/project


%% TITLE
\title{Relatório Fase 2}

%% AUTHORS
\author{%
\textbf{Maria Fernandes da Silva} \hspace{0.5cm} {\small A106802} \\
\textbf{Matilde Rodrigues Teixeira} \hspace{0.5cm} {\small A106876} \\
\textbf{Francisco Duarte Fernandes Silva} \hspace{0.5cm} {\small A100752}
}

%% Date
\date{\today}

%% Course
\newcommand{\Course}{Licenciatura de Engenharia Informática}

%% Department
\newcommand{\Department}{Escola de Engenharia}

%% UniName
\newcommand{\UniName}{Universidade do Minho}

%% UniPic
\newcommand{\UniPic}{\includegraphics[scale=0.09]{images/uminho.png}}

%% University 
\newcommand{\University}{
    \begin{flushleft}
        \UniPic
    \end{flushleft}
    \textcolor{gray}{\small\textbf{\textsf{\UniName}}}\par
    \textcolor{gray!80!white}{\small{\textsf{\Department}}}\par
    \textcolor{gray!70!white}{\small{\textsf{\Course}}}
}

%% UC
\newcommand{\UC}{
    \begin{flushleft}
        \par\textcolor{titlepagecolor}{  \LARGE\textbf{\textsf{Laboratórios de Informática III}}}
    \end{flushleft}
}

%% School Year
\newcommand{\SchoolYear}{
    \small{\textsf{Ano Letivo de 2025/2026}}}


%% Define new command to show title, author and date
\makeatletter
\let\Title\@title
\let\Author\@author
\let\Date\@date
\makeatother

%==========================================================================
% CLASSIFICATION SECTION 
%==========================================================================

%% School Year
\newcommand{\ReceptionDate}{}
%% Responsible
\newcommand{\Responsible}{}
%% Evaluation
\newcommand{\Evaluation}{}
%% Observations
\newcommand{\Observations}{}





%% MAKETEMPLATE
\newcommand{\makecover}{

%==========================================================================
% BEGIN COVER PAGE 
%==========================================================================

%% Removes page number on footer
\thispagestyle{empty}

%% No indentation 
\setlength{\parindent}{0em}

%% Put Background defined on \backgroundsetup, in this page
\BgThispage

%% Changing geometry to prevent overlay with text
%% At the end of back cover, geometry is default with \restoregeometry
\newgeometry{top=4cm,left=6cm,right=3cm,bottom=2cm}

%% builds university info defined previously
\University
\vspace{1cm}
%% builds curricular unity info defined previously
\UC
%% builds school year info defined previously
\SchoolYear

\vspace*{4cm}
%% bigger space (i think its the default one) between paragraphs 
\setlength{\parskip}{1em}

%% builds title info defined previously
\par\textbf{\textsf{\huge\Title}}
\vspace{1.5cm}
%% builds author(s) info defined previously
\par\Author

\vspace{0.5cm}

%% builds date info defined previously
\par\Date
\restoregeometry
\pagebreak

%==========================================================================
% END COVER PAGE 
%==========================================================================
}
 %% Costume made cover
\pagenumbering{gobble}
\makecover % builds the cover
\newgeometry{top=3cm,left=3cm,right=3cm,bottom=4cm} %% smaller footer and header size

%% BEGIN ABSTRACT PAGE
\renewenvironment{abstract}
 {\par\noindent\textbf{\Large\abstractname}\par\bigskip}
 {}

\pagebreak

%% BEGIN INDEXES PAGES
\renewcommand{\contentsname}{Índice}
\tableofcontents
\pagebreak


%% BEGIN INTRODUCTION
\pagenumbering{arabic}
\chapter{Introdução}

O presente relatório documenta o desenvolvimento da segunda fase do trabalho prático da unidade curricular de Laboratórios de Informática III (LI3). Este projeto teve como principal objetivo consolidar competências na linguagem C, com especial foco nos conceitos de modularidade e encapsulamento, fundamentais para o desenvolvimento de software de qualidade.

A aplicação desenvolvida consiste num sistema de gestão de informação aeroportuária, capaz de processar e responder a interrogações (queries) sobre dados de aeroportos, voos, aeronaves, passageiros e reservas. A Fase 2 representou uma evolução significativa face à primeira entrega, introduzindo três novas queries (Q4, Q5 e Q6), além de melhorias na organização do código e na gestão de memória.

O desenvolvimento seguiu uma arquitetura modular, separando claramente as responsabilidades entre camadas de apresentação, lógica de negócio, repositórios de dados e parsing. Esta organização permitiu não só uma melhor manutenção do código, mas também facilitou a implementação das novas funcionalidades exigidas.

Ao longo deste documento, será apresentada a arquitetura do sistema, as estruturas de dados utilizadas, a análise de desempenho obtida na plataforma de testes, bem como uma reflexão sobre os aspetos conseguidos e aqueles que poderiam ser melhorados.

%% ======================= %%
\chapter{Arquitetura do Sistema}

\section{Organização Geral}

A estrutura do projeto segue uma organização hierárquica que separa claramente os diferentes componentes do sistema. A divisão em pastas reflete as diferentes camadas da aplicação:

\begin{verbatim}
trabalho-pratico/
├── include/
│   ├── logica/
│   │   ├── entidades/
│   │   │   ├── aircraft.h
│   │   │   ├── airport.h
│   │   │   ├── flight.h
│   │   │   ├── passenger.h
│   │   │   └── reservation.h
│   │   ├── comando.h
│   │   ├── dados.h
│   │   ├── gestor-queries.h
│   │   ├── query1.h ... query6.h
│   │   ├── resultado-query.h
│   │   └── validacoes.h
│   ├── parsers/
│   │   ├── csv.h
│   │   └── parser.h
│   ├── repositorios/
│   │   ├── repositorio_aircrafts.h
│   │   ├── repositorio_airports.h
│   │   ├── repositorio_flights.h
│   │   ├── repositorio_passengers.h
│   │   └── repositorio_reservations.h
│   └── utils.h
├── src/
│   ├── apresentacao/
│   │   ├── programa-principal.c
│   │   └── programa-testes.c
│   ├── logica/
│   ├── parsers/
│   └── repositorios/
└── resultados/
\end{verbatim}

Esta organização traz diversos benefícios práticos: facilita a navegação no código, permite que diferentes membros do grupo trabalhem em módulos distintos sem conflitos, e torna o processo de compilação incremental mais eficiente através do Makefile.

\section{Camadas da Aplicação}

O sistema foi desenvolvido seguindo uma arquitetura em camadas, onde cada uma tem responsabilidades bem definidas:

\subsection{Camada de Apresentação}

Localizada em \texttt{src/apresentacao/}, contém os pontos de entrada da aplicação:
\begin{itemize}
    \item \texttt{programa-principal.c}: executa o carregamento dos dados e processamento das queries a partir de um ficheiro de input;
    \item \texttt{programa-testes.c}: compara os resultados obtidos com os esperados, reportando estatísticas de correção.
\end{itemize}

\subsection{Camada de Lógica}

Situada em \texttt{src/logica/}, implementa toda a lógica de negócio:
\begin{itemize}
    \item \textbf{Entidades}: estruturas de dados que representam os objetos do domínio (Airport, Flight, Aircraft, Passenger, Reservation);
    \item \textbf{Queries}: implementação das seis interrogações (Q1 a Q6);
    \item \textbf{Gestor de Queries}: coordena a execução das queries através de um dispatcher centralizado;
    \item \textbf{Validações}: funções para validar datas, emails e outros campos dos dados de entrada.
\end{itemize}

\subsection{Camada de Repositórios}

Em \texttt{src/repositorios/}, encapsula o acesso aos dados:
\begin{itemize}
    \item Cada entidade possui o seu próprio repositório;
    \item Os repositórios utilizam hash tables da GLib para acesso eficiente;
    \item Mantêm estruturas auxiliares pré-computadas para otimizar queries específicas.
\end{itemize}

\subsection{Camada de Parsing}

Localizada em \texttt{src/parsers/}, trata do carregamento dos ficheiros CSV:
\begin{itemize}
    \item \texttt{csv.c}: funções genéricas para parsing de linhas CSV;
    \item \texttt{parser.c}: coordena o carregamento de todos os ficheiros de dados.
\end{itemize}

\section{Modularidade e Encapsulamento}

Um dos objetivos centrais da cadeira de LI3 é a aplicação prática dos conceitos de modularidade e encapsulamento. No nosso projeto, estes princípios foram implementados da seguinte forma:

\subsection{Separação Interface/Implementação}

Todos os módulos seguem a convenção de separar a interface (ficheiro \texttt{.h}) da implementação (ficheiro \texttt{.c}). Os headers contêm apenas as declarações públicas necessárias, com guardas de inclusão para evitar múltiplas definições:

\begin{lstlisting}[caption={Exemplo de header com guarda de inclusão}]
#ifndef DADOS_H
#define DADOS_H

typedef struct Dados Dados;

Dados *dados_novo(void);
void   dados_destroi(Dados *dados);

#endif
\end{lstlisting}

\subsection{Estruturas Opacas}

Para garantir o encapsulamento, várias estruturas são declaradas como tipos opacos nos headers, sendo a sua definição completa visível apenas nos ficheiros de implementação. Isto impede que módulos externos acedam diretamente aos campos internos:

\begin{lstlisting}[caption={Estrutura opaca - declaração no header}]
typedef struct Dados Dados;
\end{lstlisting}

\begin{lstlisting}[caption={Estrutura opaca - definição no .c}]
struct Dados {
    RepositorioAirports     *repositorioAirports;
    RepositorioFlights      *repositorioFlights;
    RepositorioReservations *repositorioReservations;
    RepositorioAircrafts    *repositorioAircrafts;
    RepositorioPassengers   *repositorioPassengers;
    GHashTable *q6_arrivals_por_nacionalidade;
};
\end{lstlisting}

\subsection{Getters e Setters}

O acesso aos dados internos é feito exclusivamente através de funções getter, garantindo que alterações na representação interna não afetam os módulos clientes:

\begin{lstlisting}[caption={Exemplo de getter}]
RepositorioAirports *
dados_get_repositorioAirports(Dados *dados) {
    return dados ? dados->repositorioAirports : NULL;
}
\end{lstlisting}

%% ======================= %%
\chapter{Estruturas de Dados}

\section{Escolha das Estruturas}

A escolha das estruturas de dados foi orientada pelos requisitos de desempenho das diferentes queries. Optámos por utilizar as estruturas disponibilizadas pela biblioteca GLib, nomeadamente \texttt{GHashTable} e \texttt{GList}, por oferecerem implementações robustas e bem testadas.

\section{Repositórios}

Cada entidade do sistema possui um repositório dedicado que encapsula a sua gestão. A tabela seguinte resume as estruturas utilizadas:

\begin{table}[H]
\centering
\begin{tabular}{|l|l|l|}
\hline
\textbf{Entidade} & \textbf{Estrutura Principal} & \textbf{Complexidade Acesso} \\
\hline
Airports & GHashTable (by\_code) & O(1) \\
\hline
Aircrafts & GHashTable (by\_identifier) & O(1) \\
\hline
Flights & GHashTable (by\_id) + auxiliares & O(1) \\
\hline
Passengers & GHashTable (by\_document) & O(1) \\
\hline
Reservations & GHashTable (by\_id) & O(1) \\
\hline
\end{tabular}
\caption{Estruturas de dados por entidade}
\end{table}

\section{Estruturas Auxiliares para Queries}

Para otimizar o desempenho das queries, foram criadas estruturas auxiliares que pré-computam informação durante o carregamento dos dados:

\subsection{Query 2 - Aeronaves com Mais Voos}

O repositório de voos mantém uma hash table \texttt{count\_por\_aircraft} que conta o número de voos por aeronave. Esta contagem é atualizada durante a inserção de cada voo válido, evitando iterações completas durante a execução da query.

\subsection{Query 3 - Aeroporto com Mais Partidas}

A estrutura \texttt{departures\_por\_data} organiza as partidas por data e aeroporto numa hash table de dois níveis:
\begin{itemize}
    \item Primeiro nível: data (string ``YYYY-MM-DD'')
    \item Segundo nível: código do aeroporto → contagem
\end{itemize}

\subsection{Query 5 - Companhias com Maior Atraso}

A hash table \texttt{delay\_por\_companhia} acumula estatísticas de atraso por companhia aérea, incluindo o número total de voos atrasados e a soma dos minutos de atraso.

\subsection{Query 6 - Destinos por Nacionalidade}

No módulo de dados, a hash table \texttt{q6\_arrivals\_por\_nacionalidade} mapeia nacionalidades para aeroportos de chegada, permitindo responder à query sem iterar sobre todos os passageiros.

%% ======================= %%
\chapter{Implementação das Queries}

\section{Query 1 - Informação de Aeroporto}

A Q1 recebe um código IATA e devolve informação detalhada sobre o aeroporto correspondente, incluindo contagens de chegadas e partidas.

A implementação é direta: pesquisa-se o aeroporto no repositório e formata-se o resultado. A complexidade é O(1) graças à hash table.

\section{Query 2 - Top N Aeronaves}

A Q2 lista as N aeronaves com mais voos, opcionalmente filtradas por fabricante.

O algoritmo:
\begin{enumerate}
    \item Obtém a hash table de contagens por aeronave;
    \item Filtra por fabricante (se especificado);
    \item Ordena usando \texttt{qsort} com critério: contagem descendente, depois identificador alfabético;
    \item Retorna os primeiros N resultados.
\end{enumerate}

\section{Query 3 - Aeroporto com Mais Partidas}

A Q3 determina o aeroporto com mais partidas num intervalo de datas.

A implementação itera sobre as datas no intervalo especificado, acumulando as contagens por aeroporto. O aeroporto com maior total é retornado, com desempate alfabético.

\section{Query 4 - Passageiro com Mais Semanas no Top 10}

A Q4 identifica o passageiro que permaneceu mais semanas no top 10 de gastos semanais.

\textbf{Nota:} Esta query encontra-se parcialmente implementada no nosso projeto. A estrutura base está definida, mas a lógica de cálculo semanal não foi completamente desenvolvida, o que explica os resultados mais baixos obtidos na plataforma de testes.

\section{Query 5 - Companhias com Maior Atraso Médio}

A Q5 lista as N companhias aéreas com maior atraso médio, considerando apenas voos com status ``Delayed''.

A implementação utiliza as estatísticas pré-computadas no repositório de voos, ordenando por média de atraso descendente e nome alfabético em caso de empate. A precisão de três casas decimais é garantida através de formatação apropriada.

\section{Query 6 - Destino Favorito por Nacionalidade}

A Q6 identifica o aeroporto de destino mais frequente para passageiros de uma dada nacionalidade.

A normalização da nacionalidade (uppercase, remoção de espaços) garante correspondência correta independentemente do formato de entrada. A estrutura auxiliar permite resposta em tempo constante.

%% ======================= %%
\chapter{Validação de Dados}

\section{Abordagem Geral}

A validação dos dados de entrada é crucial para garantir a integridade das respostas às queries. Implementámos validações específicas para cada tipo de campo:

\subsection{Datas e Timestamps}

\begin{itemize}
    \item Formato: ``YYYY-MM-DD'' para datas, ``YYYY-MM-DD HH:MM'' para timestamps;
    \item Intervalo válido: 1900 a 30/09/2025;
    \item Valor especial ``N/A'' aceite para campos opcionais.
\end{itemize}

\subsection{Emails}

Validação de estrutura básica: presença de ``@'', caracteres antes e depois, e domínio com pelo menos um ponto.

\subsection{Campos Obrigatórios}

Verificação de preenchimento de todos os campos obrigatórios de cada entidade.

\section{Gestão de Erros}

Os registos inválidos são escritos em ficheiros de erro separados (e.g., \texttt{flights\_errors.csv}), mantendo o header original para facilitar a análise posterior.

%% ======================= %%
\chapter{Programa de Testes}

\section{Funcionamento}

O programa de testes (\texttt{programa-testes}) automatiza a validação das queries, comparando os ficheiros de output gerados com os resultados esperados.

O fluxo de execução é:
\begin{enumerate}
    \item Recebe como argumentos: diretório do dataset, ficheiro de input e diretório com outputs esperados;
    \item Executa o \texttt{programa-principal} com o dataset e input fornecidos;
    \item Compara cada \texttt{command*\_output.txt} gerado com o correspondente esperado;
    \item Reporta estatísticas por query e discrepâncias encontradas.
\end{enumerate}

\section{Métricas Reportadas}

O programa reporta:
\begin{itemize}
    \item Número de testes corretos vs. total por query;
    \item Linha da primeira discrepância em cada ficheiro com erro;
    \item Memória máxima utilizada (via \texttt{getrusage});
    \item Tempo total de execução.
\end{itemize}

%% ======================= %%
\chapter{Análise de Resultados}

\section{Metodologia}

A análise de desempenho foi realizada utilizando a plataforma oficial de testes da cadeira, que executa 300 queries (50 por cada Q1-Q6) em dois datasets distintos: regular e large.

\section{Resultados Obtidos}

\subsection{Regular Dataset}

\begin{table}[H]
\centering
\begin{tabular}{|l|c|c|c|}
\hline
\textbf{Query} & \textbf{Correção} & \textbf{Corretos/Total} \\
\hline
Q1 & 32\% & 16/50 \\
\hline
Q2 & 100\% & 50/50 \\
\hline
Q3 & 100\% & 50/50 \\
\hline
Q4 & 20\% & 10/50 \\
\hline
Q5 & 100\% & 50/50 \\
\hline
Q6 & 72\% & 36/50 \\
\hline
\end{tabular}
\caption{Resultados no Regular Dataset}
\end{table}

\textbf{Métricas gerais:}
\begin{itemize}
    \item Tempo de execução: 57.6 segundos
    \item Memória utilizada: 674 MB
    \item Memory leaks: 0.195 MB
\end{itemize}

\subsection{Large Dataset}

\begin{table}[H]
\centering
\begin{tabular}{|l|c|c|c|}
\hline
\textbf{Query} & \textbf{Correção} & \textbf{Corretos/Total} \\
\hline
Q1 & 32\% & 16/50 \\
\hline
Q2 & 100\% & 50/50 \\
\hline
Q3 & 100\% & 50/50 \\
\hline
Q4 & 30\% & 15/50 \\
\hline
Q5 & 100\% & 50/50 \\
\hline
Q6 & 0\% & 0/50 \\
\hline
\end{tabular}
\caption{Resultados no Large Dataset}
\end{table}

\textbf{Métricas gerais:}
\begin{itemize}
    \item Tempo de execução: 327.5 segundos
    \item Memória utilizada: 4502 MB
    \item Memory leaks: 0.195 MB
\end{itemize}

\section{Análise dos Resultados}

\subsection{Queries com Sucesso Completo}

As queries Q2, Q3 e Q5 obtiveram 100\% de correção em ambos os datasets, demonstrando que:
\begin{itemize}
    \item A lógica de contagem de voos por aeronave está correta;
    \item O cálculo de partidas por aeroporto e intervalo funciona adequadamente;
    \item As estatísticas de atraso por companhia são calculadas corretamente.
\end{itemize}

\subsection{Queries com Problemas Identificados}

\textbf{Q1 (32\%):} A baixa taxa de sucesso sugere problemas na contagem de chegadas/partidas ou na validação de aeroportos. Possíveis causas:
\begin{itemize}
    \item Contagem incorreta de passageiros por aeroporto;
    \item Problemas na atualização dos contadores durante o parsing de reservas.
\end{itemize}

\textbf{Q4 (20-30\%):} Como mencionado anteriormente, esta query não foi completamente implementada. O código presente é estrutural, mas falta a lógica de cálculo semanal e ranking.

\textbf{Q6 (72\% regular, 0\% large):} A discrepância entre datasets indica um problema de escalabilidade ou sensibilidade a casos específicos do large dataset. Possíveis causas:
\begin{itemize}
    \item Normalização incompleta de nacionalidades;
    \item Problemas no mapeamento passageiro-voo-aeroporto;
    \item Casos especiais não tratados no dataset maior.
\end{itemize}

\subsection{Gestão de Memória}

O valor de memory leaks (0.195 MB) indica uma boa gestão de memória, com apenas pequenas fugas residuais. O consumo de memória escala de forma aproximadamente linear com o tamanho do dataset (674 MB → 4502 MB), o que é expectável dada a natureza das estruturas utilizadas.

\subsection{Desempenho Temporal}

O tempo de execução aumenta significativamente no large dataset (57.6s → 327.5s), um fator de aproximadamente 5.7x. Este aumento é superior ao esperado para um crescimento linear, sugerindo que algumas operações têm complexidade superior a O(n).

%% ======================= %%
\chapter{Aspetos a Melhorar}

\section{Correção das Queries}

\subsection{Query 1}

É necessário rever a lógica de contagem de chegadas e partidas. Suspeita-se que o problema esteja na forma como as reservas atualizam os contadores dos aeroportos.

\subsection{Query 4}

Esta query requer uma reimplementação completa. O algoritmo deve:
\begin{enumerate}
    \item Calcular os gastos semanais de cada passageiro;
    \item Determinar o top 10 para cada semana;
    \item Contar as ocorrências de cada passageiro no top 10;
    \item Ordenar e devolver o resultado.
\end{enumerate}

\subsection{Query 6}

Investigar a diferença de comportamento entre datasets. Recomenda-se adicionar logs de debug para identificar onde ocorre a divergência.

\section{Otimização de Desempenho}

\subsection{Parsing}

Apesar de já utilizarmos buffers grandes (1 MB), poderiam ser exploradas técnicas adicionais como:
\begin{itemize}
    \item Memory-mapped files para datasets muito grandes;
    \item Parsing paralelo dos diferentes ficheiros CSV.
\end{itemize}

\subsection{Estruturas de Dados}

Para datasets ainda maiores, poderia ser benéfico:
\begin{itemize}
    \item Utilizar árvores balanceadas para queries que requerem ordenação frequente;
    \item Implementar índices secundários para acelerar pesquisas por múltiplos critérios.
\end{itemize}

\section{Qualidade do Código}

\subsection{Remoção de Código de Debug}

O código atual contém várias instruções de debug (printf) que deveriam ser removidas ou condicionadas a uma flag de compilação.

\subsection{Documentação}

Alguns módulos carecem de documentação adequada. Recomenda-se a utilização de comentários no estilo Doxygen para facilitar a geração de documentação automática.

%% ======================= %%
\chapter{Conclusão}

\section{Objetivos Alcançados}

O desenvolvimento da Fase 2 permitiu consolidar vários conceitos fundamentais da unidade curricular:

\begin{itemize}
    \item \textbf{Modularidade}: o código está organizado em módulos coesos com responsabilidades bem definidas;
    \item \textbf{Encapsulamento}: a utilização de estruturas opacas e getters protege a implementação interna;
    \item \textbf{Gestão de memória}: o baixo nível de memory leaks demonstra atenção à libertação de recursos;
    \item \textbf{Utilização de ferramentas}: Makefile funcional, utilização da GLib para estruturas de dados.
\end{itemize}

Em termos de funcionalidade, as queries Q2, Q3 e Q5 funcionam corretamente, representando 50\% das interrogações com 100\% de sucesso.

\section{Principais Aprendizagens}

Este projeto proporcionou aprendizagens valiosas:

\begin{itemize}
    \item A importância de planear a arquitetura antes de começar a implementação;
    \item As vantagens de utilizar bibliotecas testadas (GLib) em vez de implementar estruturas de raiz;
    \item A necessidade de testes incrementais durante o desenvolvimento;
    \item Os desafios de escalar código para datasets de maior dimensão.
\end{itemize}

\section{Reflexão Crítica}

Reconhecemos que o projeto apresenta limitações, nomeadamente nas queries Q1, Q4 e Q6. Em retrospetiva, deveríamos ter:

\begin{itemize}
    \item Começado a implementação da Q4 mais cedo, dado a sua complexidade;
    \item Testado com mais frequência contra os datasets oficiais;
    \item Dedicado mais tempo à análise dos outputs incorretos para identificar padrões de erro.
\end{itemize}

Não obstante, consideramos que o projeto cumpre os objetivos pedagógicos principais da cadeira, demonstrando capacidade de estruturar um programa C de média dimensão seguindo boas práticas de engenharia de software.

\end{document}
